\documentclass[11pt]{article}
\usepackage[utf8]{inputenc}
\usepackage[T1]{fontenc}
\usepackage[margin=1in]{geometry}
\usepackage{hyperref}
\usepackage{enumitem}
\title{Magic School | Day 2}
\author{Piter Garcia}
\date{March 20, 2026}
\begin{document}
\maketitle
\section*{Quick Scan}
\begin{itemize}[leftmargin=*]
\item Grade: Grade 4
\item Workbook sessions: Session 20
\item Class blocks: Day 1 1:15-2:10 Baris, Day 2 1:15-2:10 Fowler, Day 3 1:15-2:10 Schrank, Day 4 1:15-2:10 Autore
\item Reference file: magic-school-day-2.bib
\end{itemize}
\section*{School Planning Goals and Inclusion}
\begin{itemize}[leftmargin=*]
\item What is expected: - I can explain what MagicSchool is and why it's useful for learning. | - I can log into MagicSchool and find my way around the main areas. | - I can name at least three different types of activities available on MagicSchool. | - I can use MagicSchool safely and responsibly. | - I can share one way MagicSchool helps me learn better. | - I can model how computer hardware and software work together as a system to accomplish tasks. | - I can determine potential solutions to solve hardware and software problems using common troubleshooting strategies.
\item What we achieved: This lesson points to local transcript or evidence files in the workspace so the packet can be revised from what actually happened in class.
\item What needs improvement: stronger proof, cleaner assessment notes, and tighter lesson artifacts.
\item How to improve: add transcript proof, one saved artifact, and clearer success criteria.
\item Inclusion focus: use visual modeling, chunked directions, predictable routines, and flexible participation roles.
\end{itemize}
\section*{Official References}
\begin{itemize}[leftmargin=*]
\item \url{https://csteachers.org/k12standards/interactive/}

\end{itemize}
\end{document}
